% STATUS: draft written, session 1. Revisit last, once results are final.
% LENGTH: 150-200 words.
Automatic summarization of product reviews is usually framed as a modelling
problem, but on real e-commerce data it is equally a problem of \emph{form}.
We fine-tune \texttt{t5-small}, a 60M-parameter encoder--decoder, on 4{,}000
Amazon food reviews whose reference summaries are the headlines customers
wrote themselves, and evaluate on 300 held-out reviews. Fine-tuning improves
ROUGE-1 from 0.109 to 0.175, a relative gain of 61\%. The more instructive
result is the baseline: the \emph{same pretrained model without fine-tuning}
scores 0.109, below the 0.126 of a trivial extractive baseline that simply
returns the first sentence. The un-tuned model is fluent and on-topic, but
writes full sentences where the task rewards a four-word headline; fluency is
not task fit. We further show that per-review summarization, while accurate,
systematically conceals aggregate signal: for a baby-formula product with a
3.13 star average and predominantly positive summaries, aspect mining over the
same reviews surfaces sustained complaints about arsenic levels and infant
health. We argue that review summarization systems should be evaluated at the
aggregate level, not only per review.
